% ==============================================================================
% 10_ha_tang_devops.tex — Phần 10. Thiết kế hạ tầng và vận hành (DevOps)
% Phụ trách: Xuân (Dev 2) — phối hợp Tiến phần Cloud/Cloudflare Pages
% Ngân sách trang: 3 trang
% ==============================================================================

\pttkPart{Thiết kế hạ tầng và vận hành (DevOps)}

\todosection{Xuân (Dev 2)}{Tham chiếu \texttt{docker-compose.yml}, \texttt{.github/workflows/ci.yml}, \texttt{infra/} và \texttt{docs/xuan/18..19}. Ngân sách mục tiêu: 3 trang.}

\subsection{Sơ đồ triển khai}
\todofig{Hình 10. Sơ đồ triển khai DevFlow (Local Docker Compose và Production Cloud) — diagrams/out/hinh\_10\_deployment.pdf}

\subsection{Các môi trường vận hành}
% Bảng môi trường: Local/Dev (Docker Compose), CI/Test (H2/Testcontainers), Production (Cloudflare Pages + PaaS + Managed PostgreSQL).

\subsection{Quy trình CI/CD}
% Quy trình PR -> ./gradlew :verifyModuleBoundaries -> ./gradlew check -> npm run build/lint -> Docker build -> Deploy Cloud.

\subsection{Giám sát và mục tiêu mức dịch vụ (SLO)}
% Bảng SLI/SLO (API success >= 99.5\%, p95 < 200ms, AI latency, Webhook delivery) + Actuator/Micrometer -> Prometheus -> Grafana + JSON logs.

\subsection{Sao lưu và khôi phục thảm họa}
% Bảng chiến lược sao lưu PostgreSQL PITR, RPO <= 15 phút, RTO <= 4 giờ, quản lý schema qua Flyway.
